\documentclass[11pt]{article}
\usepackage[margin=1in]{geometry}
\usepackage{booktabs}
\usepackage{hyperref}
\usepackage{amsmath}

\title{Tau Scaling as an Evidence-Gated Public Claim System\\
\large Source Population, Primary-Validation Gaps, and Non-Claim Locks}
\author{James Paul Jackson}
\date{TAU-SCALING-SA v0.9.1}

\begin{document}
\maketitle

\begin{abstract}
This manuscript scaffold presents Tau Scaling as a public semiconductor-claim governance problem rather than as a validated silicon result. The repository evaluates public Tau Scaling and LogicFolding claims through a reproducible evidence pipeline. In the current repository state, eight public claims are source-populated from public reporting, zero are primary-source confirmed, and zero are independently confirmed. The contribution is therefore a reproducible source-provenance and validation-gap artifact, not a claim of silicon validation, product validation, benchmark superiority, process-node equivalence, or a universal Tau Scaling law.
\end{abstract}

\section{Introduction}
Public semiconductor claims often blend roadmap language, media interpretation, methodology framing, reported metrics, and implied validation. The tau-scaling repository treats this as a claim-governance problem.

\section{Method}
The method converts public claims into a staged evidence pipeline:
claim ledger, source provenance, source intake cards, public source population, primary-source validation gap review, milestone package, and manuscript scaffold.

\section{Results}
\begin{center}
\begin{tabular}{lr}
\toprule
Metric & Value\\
\midrule
Public claims analyzed & 8\\
Source-populated claims & 8\\
Primary-source confirmed claims & 0\\
Independent-source confirmed claims & 0\\
High validation-gap claims & 8\\
Release findings & 0\\
\bottomrule
\end{tabular}
\end{center}

\section{Interpretation}
The current evidence state supports a bounded claim: Tau Scaling can be represented as an evidence-gated public claim system. It does not validate silicon or promote any technical claim beyond available evidence.

\section{Limitations}
This artifact does not validate semiconductor performance, Huawei product performance, manufacturing capability, process-node equivalence, benchmark superiority, or a universal Tau Scaling law.

\section{Reproducibility}
All source-provenance, validation-gap, release-readiness, and manuscript-scaffold artifacts are stored in the repository under reports, releases, sources, visuals, and docs.

\section{Non-Claim Locks}
Public reporting is not primary validation. Source population is not source validation. Source validation is not claim promotion. A roadmap target is not an achieved result.

\end{document}
